\scanpage{101}
\mspage{96}

\begin{center}
Chapter III

Some Applications of the Criterion of Projectivity

of Abstract Complete Varieties.
\end{center}

In this chapter we will examine some of the applications of the following assertion which follows from theorem 2:

\uline{Let} $V$ \uline{be a nonsingular complete algebraic variety} $\dim V=s$ \uline{and} $D$ \uline{be a divisor on} $V$ \uline{such that} $D^kW_k>0$ \uline{for any irreducible subvariety} $W_k$ \uline{of dimension} $k$, $k=\overline{1,s}$ ($D^sV=D^s$). \uline{Then, at sufficiently large} $n$, \uline{the mapping of the variety} $V$ \uline{into a projective space, corresponding to the fibre} $[nD]$, \uline{becomes a biregular imbedding}.

Let $V$ be a complete normal algebraic variety, $\dim V=s$ and $U$ be an open set on $V$, being a quasi-projective variety (that is, allowing realization in the form of an open subset of a projective algebraic variety).
\margin{p. 218}
Let $x_0,x_1,\ldots$ be homogeneous coordinates, corresponding to $U$, $x_0\ne0$, and $F_j=\dfrac{x_j}{x_0}$. Let $\unsure{\{F\}|[F]}$ be \deleted{a}\inserted{\unsure{the}} linear space over $k$, generated by the functions $F_j$, and $\{T_{\inserted{\ell}}\}$ be a totality of irreducible $s-1$ dimensional components of the set $V-U$,
\margin{$\ell$}
\[
\tau_\ell=\min_j n_{T_\ell}(F_j)
\]
where $n_{T_l}(F_j)$ denotes the order of the function $F_j$ on the divisor $T_l$. Let a function $\varphi\in\unsure{\{F\}|[F]}$ be such that,

\margin{100}
